\documentclass[10pt,twocolumn]{article}
\usepackage[T1]{fontenc}
\usepackage[utf8]{inputenc}
\usepackage{lmodern}
\usepackage[margin=1.8cm,top=2.4cm,bottom=2cm,columnsep=0.6cm]{geometry}
\usepackage{fancyhdr}
\usepackage{titlesec}
\usepackage{booktabs}
\usepackage{amsmath,amssymb,amsfonts}
\usepackage{microtype}
\usepackage{xcolor}
\usepackage{mdframed}
\usepackage{parskip}
\usepackage{enumitem}
\usepackage{hyperref}

\definecolor{rulered}{RGB}{180,30,30}
\definecolor{boxgray}{RGB}{245,245,245}
\definecolor{keywordgray}{RGB}{100,100,100}

\pagestyle{fancy}
\fancyhf{}
\fancyhead[L]{\textsc{\small Claude Mag}}
\fancyhead[C]{\small\textit{Machine Learning Research Digest}}
\fancyhead[R]{\small 2026-10-26}
\fancyfoot[C]{\small\thepage}
\renewcommand{\headrulewidth}{0.8pt}

\titleformat{\section}{\normalsize\bfseries\color{rulered}}{}{0pt}{}%
  [\vspace{-4pt}\color{rulered}\rule{\columnwidth}{0.4pt}]
\titlespacing{\section}{0pt}{8pt}{4pt}

\hypersetup{colorlinks=true,urlcolor=rulered,linkcolor=black}

\begin{document}
\setlength{\parindent}{0pt}
\setlength{\parskip}{4pt}

{\small\color{keywordgray}\textbf{Keywords:}
  speculative decoding \textbullet{}
  LLM inference acceleration \textbullet{}
  online distillation \textbullet{}
  chain-of-thought}\par\vspace{4pt}

{\LARGE\bfseries\raggedright
  Test-Time Speculation}\par\vspace{4pt}

{\large\itshape\raggedright
  Continuously Adapting the Draft Model During Generation Sustains
  Speculative Speedups Throughout Long Chain-of-Thought Outputs}\par\vspace{6pt}

{\small\textbf{By} Avinash Kumar, Sujay Sanghavi, Poulami Das
  (UT Austin)}\par\vspace{2pt}

{\small\color{keywordgray}arXiv:2605.09329 \textbullet{}
  NeurIPS 2026}
\par\vspace{2pt}\noindent{\color{rulered}\rule{\columnwidth}{0.6pt}}\vspace{6pt}

Speculative decoding accelerates LLM inference by having a small draft
model propose candidate tokens that a large target model verifies in
parallel. In practice, however, the acceptance length of trained
speculators---including state-of-the-art systems EAGLE-3, DFlash, and
PARD---decays to near 1 after a few thousand output tokens, eliminating
the speedup for long chain-of-thought generation. Test-Time Speculation
(TTS) fixes this by adapting the draft model online, using the
target-model evaluations that the verification step already performs, at
zero added inference cost. On Qwen-3, Qwen-3.5, and Llama~3.1 model
families, TTS improves acceptance length by up to 72\% and 41\% on average
over the best static speculator, translating to $2.1$--$3.4\times$ wall-clock
speedup on 5K-token generation tasks.

\begin{mdframed}[backgroundcolor=boxgray,linecolor=rulered,linewidth=1pt,
    innerleftmargin=6pt,innerrightmargin=6pt,
    innertopmargin=6pt,innerbottommargin=6pt]
\textbf{\small AT A GLANCE}\par\vspace{4pt}
\begin{description}[leftmargin=0pt,labelsep=4pt,itemsep=2pt,parsep=0pt,
    font=\small\bfseries]
  \item[Problem:] Trained speculators are distilled offline on short sequences;
    their acceptance length decays toward 1 on the long CoT outputs that matter
    most for reasoning tasks.
  \item[Why it matters:] Speculative decoding is the leading inference
    acceleration technique, but its speedup evaporates precisely for the
    long-horizon tasks that consume the most compute.
  \item[Method:] Online KL distillation of the draft model using
    target-model token probabilities already computed by the verification step,
    plus a proximity regularizer to prevent forgetting.
  \item[Result:] Up to 72\% improvement in acceptance length; $2.1$--$3.4\times$
    wall-clock speedup on 5K+ token outputs.
  \item[Evidence:] MATH-500, HumanEval+, long-document summarization;
    Qwen-3/3.5 and Llama~3.1 families; ablations on regularizer and step count.
\end{description}
\end{mdframed}

\section{The Big Picture}

Speculative decoding exploits the asymmetry between generation and
verification: a small, fast draft model proposes $k$ tokens, and the large
target model verifies all $k$ in one forward pass---accepting the longest
prefix consistent with its distribution. When the draft model agrees with the
target model, the effective throughput multiplies by $k$; when it does not,
the cost is at most one wasted draft pass.

The technique's Achilles' heel is draft quality at long context lengths.
All existing speculators are trained on datasets of typical generation lengths
(hundreds to low thousands of tokens); at inference, reasoning-intensive tasks
such as MATH, LiveCodeBench, and multi-step QA generate five to ten thousand
tokens. The draft model's internal representations, tuned for short contexts,
drift from the target's as the context grows, and acceptance length collapses.

Test-Time Speculation reframes speculative decoding as a continual learning
problem: the draft model is a student that should track the target's
distribution throughout generation, not just at the start.

\section{Why Existing Approaches Fall Short}

\textbf{Static draft models (EAGLE-3, DFlash, PARD):}
All trained offline on sequences far shorter than inference-time CoT outputs.
No mechanism exists to update the draft as the context lengthens.

\textbf{Longer training sequences:}
Simply including longer sequences in offline training is expensive, requires
representative CoT data that may not be available, and still does not adapt to
the specific prompt or intermediate reasoning steps of a given inference call.

\textbf{Online re-training approaches:}
Prior online adaptation methods for NLP require periodic re-training runs with
additional target-model calls, breaking the generation pipeline and negating
the speedup.

\section{Core Mechanism}

At each speculative decoding step, the target model already evaluates the
$k$ draft tokens and produces per-token probability distributions
$\pi^*(\cdot \mid x_{<j})$ for $j = t, \ldots, t+k$. These distributions are
the exact student-teacher signal needed to distill the draft model for the
current context---and conventional speculative decoding discards them after
acceptance/rejection decisions.

TTS re-uses these distributions to take one gradient step on the draft model
after each verification pass. The gradient step minimizes a KL distillation
loss plus a proximity regularizer that prevents the draft from forgetting its
base distribution. Because the gradient is taken on a small draft model
(1--3B parameters) and no additional target-model calls are needed, the
per-step overhead is negligible.

\section{Technical Formulation}

Let $\pi^*$ denote the target model and $\pi_\theta$ the draft model with
parameters $\theta$. At step $t$, the token-level KL distillation loss over
the $k$ verified positions is:
\[
\mathcal{L}_t(\theta)
  = \frac{1}{k}\sum_{j=t}^{t+k} D_{\text{KL}}\!\left(
    \pi^*(\cdot \mid x_{<j}) \;\|\; \pi_\theta(\cdot \mid x_{<j})
  \right)
\]
The proximity regularizer prevents catastrophic forgetting:
\[
\Omega(\theta, \theta_{t-1}) = \|\theta - \theta_{t-1}\|_2^2
\]
The combined update is a single gradient step:
\[
\theta_t = \theta_{t-1} - \eta\,\nabla_\theta\!\left[
  \mathcal{L}_t(\theta_{t-1}) + \beta\,\Omega(\theta_{t-1}, \theta_{t-2})
\right]
\]
with learning rate $\eta$ following a cosine decay schedule over the
generation horizon, and regularization weight $\beta$ set so that total
parameter shift remains within a small budget.

\section{Learning or Inference Procedure}

\textbf{Standard speculative decoding (baseline):}
Frozen draft model; acceptance length decays with context length.

\textbf{TTS (proposed):}
\begin{enumerate}[itemsep=2pt,parsep=0pt]
  \item Initialize $\pi_\theta$ from any pre-trained speculator.
  \item At each verification step:
    \begin{enumerate}[label=(\alph*),itemsep=1pt,parsep=0pt]
      \item Accept/reject draft tokens as usual.
      \item Compute $\mathcal{L}_t$ using the target distributions already
        obtained.
      \item Take one gradient step to update $\theta$.
    \end{enumerate}
  \item Continue generation with the updated draft model.
\end{enumerate}

No additional target-model calls. The gradient step on the small draft model
adds $< 5\%$ latency overhead in practice.

\section{What the Guarantee Says}

The distillation loss minimized at each step is the expected log-probability
ratio between the target and draft at the \emph{current} context. Minimizing
it pushes the draft toward the target's distribution at the positions just
verified. The proximity regularizer bounds the per-step parameter shift,
ensuring that adaptation at step $t$ does not degrade draft quality at step
$t-1$'s context distribution beyond a controlled amount. Together, these
properties imply that the acceptance rate converges to a level bounded by the
target's entropy at the current context---the theoretical maximum for any
draft model---rather than decaying toward 1.

\section{Experimental Findings}

\begin{itemize}[itemsep=2pt,parsep=0pt]
  \item \textbf{MATH-500 (long mathematical reasoning):} TTS achieves $3.4\times$
    wall-clock speedup over the target at 5K tokens; the best static speculator
    achieves $1.5\times$ at the same length.
  \item \textbf{HumanEval+ (long code generation):} $2.7\times$ vs.\ $1.8\times$
    for TTS and the best static speculator respectively.
  \item \textbf{Acceptance length trajectory:} Static speculators show a monotonic
    decay from $\sim 4.5$ tokens at 500 tokens to $\sim 1.2$ at 5K tokens;
    TTS maintains $\sim 4.0$ tokens throughout.
  \item \textbf{Model families:} Gains are consistent across Qwen-3-8B,
    Qwen-3.5-32B, and Llama-3.1-70B with their respective draft models.
  \item \textbf{Compatibility:} TTS improves acceptance length for all three
    base speculators (EAGLE-3, DFlash, PARD); EAGLE-3 + TTS achieves the
    best absolute acceptance length.
\end{itemize}

\section{Ablations and Interpretation}

\textbf{No proximity regularizer:}
The draft overfits to the current context window and loses generalization to
the next step; acceptance length drops by 18\% compared to full TTS.

\textbf{Multiple gradient steps per verification:}
Two steps per round adds compute without improving acceptance; one step is
both necessary and sufficient.

\textbf{Fixed vs.\ cosine learning rate:}
Cosine decay improves stability and final acceptance length by $\sim 5\%$;
a fixed high rate causes oscillation in later generation steps.

\textbf{Gains grow with output length:}
Improvement over the base speculator is near-zero at $< 500$ tokens (where
the training distribution covers the context) and reaches its maximum at
$> 5{,}000$ tokens, confirming that TTS specifically addresses the
out-of-distribution long-context problem.

\section*{Reference}

Avinash Kumar, Sujay Sanghavi, Poulami Das.
\textbf{Test-Time Speculation}.
arXiv:2605.09329. NeurIPS 2026.
\url{https://arxiv.org/abs/2605.09329}

\end{document}
